\chapter{Sufficiency}

%---------------------------------------------------------------
\section{Context and Motivation}
%---------------------------------------------------------------

Global climate targets are under pressure. The 6th IPCC report (Working Group 3)
confirms with high confidence that the world is behind on addressing climate change.
Most 2°C scenarios (two thirds of them) rely on negative carbon technologies such as
Direct Air Capture (DAC) or Bioenergy with Carbon Capture and Storage (BECCS), yet
these technologies have been studied for over 50 years without becoming economically
viable.

The remaining global carbon budgets are tight. Updated 2023 estimates put the budget
at approximately 260~GtCO\textsubscript{2} for a 50\% chance of limiting warming to
1.5°C, and around 950~GtCO\textsubscript{2} for a 67\% chance of staying below 2°C.
For context, proven fossil fuel reserves embed some 3\,500~GtCO\textsubscript{2} of
potential emissions --- far beyond any safe budget.

%---------------------------------------------------------------
\section{The Sufficiency--Efficiency--Renewables Framework}
%---------------------------------------------------------------

The presentation applies the \emph{negaWatt} approach, built on three complementary
levers formalized through the Kaya identity:

\begin{itemize}[noitemsep]
    \item \textbf{Energy sufficiency} --- reducing the level of a need through
          behavioural and societal change (e.g.\ smaller dwellings, less car use,
          dietary shifts). Introduced in IPCC AR6 WG3 as a policy category in its
          own right.
    \item \textbf{Energy efficiency} --- reducing the energy required to deliver a
          given service through better technology (heat pumps, insulation,
          electric vehicles).
    \item \textbf{Renewable energy} --- decarbonizing the remaining energy supply.
\end{itemize}

A key conceptual distinction is made between sufficiency and efficiency: efficiency
improvements can be cancelled out by growing consumption (the rebound effect). German
residential data illustrate this clearly --- gains in energy intensity per m² have
been largely offset by rising living space per capita. The Kaya identity is decomposed as:
\[
\text{Global Emissions} = \text{Population} \times \frac{\text{Prod \& Services}}{\text{Population}} \times \frac{\text{Energy}}{\text{Prod \& Services}} \times \frac{\text{Emissions}}{\text{Energy}} - \text{Negative Emissions}
\]
So there are 3 levers to reduce emmissions : Energy sufficiency (human behaviour change), Energy efficiency and Renewable energy (technical solutions) and negative emissions (carbon capture and storage, carbon capture and utilization, direct air capture, land use and forestry) which are uncertain solutions.

%---------------------------------------------------------------
\section{Sectoral Demand Scenarios}
%---------------------------------------------------------------

The negaWatt Belgium scenario covers agriculture, residential, tertiary, mobility,
freight, industry, non-energy-demand (NED) industry, and aviation bunkers.

\subsection{Buildings}

Residential space heating (65~TWh in 2019) is projected to fall by 71\% to
19~TWh in 2050, driven by:

\begin{itemize}[noitemsep]
    \item Temperature setpoint reduction: $21\,^\circ\text{C} \to 19\,^\circ\text{C}$ ($-12\%$)
    \item Building envelope renovation at 3\%/year from 2030 ($-28\%$)
    \item Heating technology switch (18\% $\to$ 82\% heat pumps and district heating) ($-20\%$)
    \item Improved system efficiency, average COP $2.5 \to 3.5$ ($-10\%$)
\end{itemize}

Combined with the tertiary sector, total building final energy consumption falls by
48\% between 2019 and 2050, with natural gas and liquid fuels being largely replaced
by ambient heat and electricity.

\subsection{Transport}

Individual car mobility (56~TWh in 2019) is reduced by 93\% to just 4~TWh by 2050
through a combination of sufficiency and efficiency measures:

\begin{itemize}[noitemsep]
    \item Distance reduction ($-10\%$) via teleworking and urban densification
    \item Modal shift: car share $80\% \to 57\%$ ($-14\%$)
    \item Car occupancy: $1.25 \to 2$ persons per vehicle ($-18\%$)
    \item Near-full fleet electrification, 99.5\% by 2050 ($-31\%$)
    \item Vehicle energy efficiency gains: $0.22 \to 0.13$~kWh/km ($-19\%$)
\end{itemize}

The number of cars on Belgian roads would fall from 6 million to 2 million.
Domestic transport final energy consumption declines by 83\% overall.

\subsection{Industry}

Key industrial transformations include:
\begin{itemize}[noitemsep]
    \item \textbf{Steel}: $-15\%$ volume reduction; fuel switch from coke to
          hydrogen (60\% direct reduction) and electric arc furnaces (40\%);
          energy intensity $3\,950 \to 2\,300$~MWh/kt.
    \item \textbf{Cement}: $-25\%$ volume (more timber construction); energy
          intensity $904 \to 650$~MWh/kt; fuel switch coal $\to$ biomass.
    \item \textbf{Chemicals}: $-20\%$ demand; process electrification and
          fuel switching to renewables.
\end{itemize}

Total industrial final energy consumption falls by 34\%.

\subsection{Overall Final Energy Consumption}

Belgium's total final energy consumption decreases from \textbf{490~TWh in 2019}
to \textbf{246~TWh in 2050}, a reduction of \textbf{50\%}, broken down as:
buildings $-49\%$, transport $-76\%$, industry $-34\%$.

%---------------------------------------------------------------
\section{Energy System Modelling: PyPSA-Belgium}
%---------------------------------------------------------------

The supply side is modelled using the open-source PyPSA-Eur framework adapted for
Belgium. The model features:

\begin{itemize}[noitemsep]
    \item One node per country, covering five countries to represent
          cross-border interconnections
    \item Optimization runs for 2030, 2040, and 2050 to map the full
          transition pathway
    \item Hourly time resolution to capture renewable variability and
          flexibility requirements
    \item Full sector coupling: electricity, gas, hydrogen, heating, and industry
\end{itemize}

Sufficiency demand assumptions from the negaWatt Belgium scenario are fed directly
as inputs to the supply-side optimization.

%---------------------------------------------------------------
\section{Two Scenarios Compared}
%---------------------------------------------------------------

\begin{center}
\begin{tabular}{lcc}
\toprule
 & \textbf{Reference (Ref)} & \textbf{Sufficiency (nW BE)} \\
\midrule
Emission target (Belgium, 2050) & $-95\%$ & $-95\%$ \\
Target (5 countries, 2050) & Net zero & Net zero \\
Demand & Current trends & Sufficient demand \\
Carbon Capture and Storage (CCS) & Allowed & Not allowed \\
Carbon Capture and Utilization (CCU) & Allowed & Allowed \\
Direct Air Capture (DAC) & Allowed & Not allowed \\
Building renovation rate & 1\%/year & 3\%/year \\
Industry & Continued trends & Increased circularity \\
\bottomrule
\end{tabular}
\end{center}

%---------------------------------------------------------------
\section{Key Results}
%---------------------------------------------------------------

\subsection{Cost-Effectiveness}

The negaWatt Belgium scenario reduces total average annual system costs
(CAPEX + OPEX) by \textbf{33\% compared to the reference scenario} and by 19\%
compared to 2019 levels. Savings arise from reduced fuel import costs, avoided
infrastructure investment, and lower operating costs due to a smaller overall
energy system.

\subsection{Energy Independence}

The sufficiency scenario doubles Belgium's share of local energy production
(from approximately 50\% to 62\%) and reduces energy imports by a factor of four.
This significantly improves strategic autonomy, especially in combination with
circular economy policies that limit dependence on critical raw materials.

\subsection{Renewable Deployment}

Sufficiency substantially reduces the renewable capacity required:

\begin{itemize}[noitemsep]
    \item \textbf{Solar PV}: 48~GW (reference) vs.\ 48~GW (sufficiency), but
          annual installation rate drops from 3\,000~MW/year to 1\,600~MW/year.
    \item \textbf{Onshore wind}: 14~GW (reference) vs.\ 7.4~GW (sufficiency);
          annual additions 440~MW vs.\ 160~MW.
    \item \textbf{Offshore wind}: 8~GW in both scenarios.
\end{itemize}

\subsection{Hydrogen}

Hydrogen demand in 2050 reaches 82~TWh in the sufficiency scenario, which is
\textbf{43\% less} than the 144~TWh needed in the reference case. Hydrogen is
reserved for hard-to-decarbonize uses: industry, long-distance transport,
synthetic fuel production, and non-energy applications. Most supply (84\%) is
imported; local electrolysis provides 13~TWh.

\subsection{Cumulative Emissions}

The sufficiency scenario achieves the $-95\%$ emission target for Belgium by 2050
without relying on CCS or DAC, using instead biomass-based carbon cycles, land use,
and forestry (LULUCF) sinks. The reference scenario reaches the same 2050 target
but relies heavily on large-scale carbon removal, resulting in significantly higher
cumulative emissions over the transition period.

%---------------------------------------------------------------
\section{Sensitivity Analyses}
%---------------------------------------------------------------

\subsection{Additional Offshore Wind}

Expanding offshore capacity by an additional 16~GW (to 24~GW total by 2050,
consistent with the Esbjerg and Ostend declarations) reduces onshore wind needs
by 35\% and increases power-to-gas capacity by 66\%. Electricity and hydrogen
imports also decrease. Core messages on cost-effectiveness and feasibility remain
unchanged.

\subsection{Nuclear Energy Costs}

At the default cost assumption of 8\,594~EUR/kW, new nuclear is not selected by
the optimizer. Nuclear begins to appear in the solution only at approximately
4\,500~EUR/kW --- a dramatic cost reduction only achievable with small modular
reactors (SMRs) under optimistic assumptions. At that price point, around 10~GW
of nuclear capacity would be deployed, substituting for some onshore wind,
solar PV, and reducing transmission needs.

%---------------------------------------------------------------
\section{Key Takeaways}
%---------------------------------------------------------------

\begin{enumerate}
    \item \textbf{Cheapest scenario}: The negaWatt Belgium pathway is
          \emph{more} cost-effective than a technology-only reference scenario,
          reducing system costs by 33\%.
    \item \textbf{Strategic independence}: Consuming less energy directly
          increases Belgium's energy autonomy and reduces exposure to
          geopolitical energy risks.
    \item \textbf{Realistic scenario}: The sufficiency scenario avoids CCS
          and DAC, technologies that face major industrial-scale deployment
          challenges. Renewable deployment remains ambitious but stays well
          below estimated technical potentials.
    \item \textbf{Fair scenario}: The scenario prioritizes collective solutions
          benefiting the most vulnerable citizens, treats affordable energy
          access as a fundamental right, and frees up global resources for
          the Global South by reducing Belgium's overall footprint.
\end{enumerate}

The full model and data are openly available at:
\begin{itemize}[noitemsep]
    \item \url{https://github.com/IntSusEnergySystems/PyPSA-Belgium}
    \item \url{https://www.negawatt.be/scenario/}
\end{itemize}


